\section{Modelo elétrico da bobina distribuída}

Em baixas frequências, uma bobina pode ser aproximada por uma indutância e uma resistência equivalentes. Em um surto rápido, essa simplificação é limitada: a capacitância do enrolamento para terra e a distribuição espacial de indutância e resistência passam a determinar como a tensão se reparte ao longo da bobina. Por isso o projeto usa uma aproximação de parâmetros distribuídos, semelhante ao tratamento de linhas por seções discretas \cite{greenwood1991}.

No caso base, os parâmetros totais são divididos em \Nsec{} seções:

\begin{align}
R_{\mathrm{sec}} &= \frac{R_{\mathrm{total}}}{N} = \frac{5}{20} = \SI{0.25}{\ohm},\\
L_{\mathrm{sec}} &= \frac{L_{\mathrm{total}}}{N} = \frac{0.01}{20} = \SI{0.5}{mH},\\
C_{\mathrm{sec}} &= \frac{C_{\mathrm{total}}}{N} = \frac{10^{-9}}{20} = \SI{50}{pF}.
\end{align}

Na topologia Pi, os capacitores de extremidade aparecem como \(\Csec/2\), enquanto os nós internos possuem \(\Csec\), resultado da soma das duas meias capacitâncias das seções vizinhas. A Figura~\ref{fig:circuito-pi} mostra essa representação.

\begin{figure}[H]
\centering
\includegraphics[width=\linewidth]{circuito_grounded.png}
\caption{Escada equivalente do enrolamento. Cada seção tem um ramo série \(R\text{--}L\), uma
capacitância série entre espiras \(\Cs\) (\emph{turn-to-turn}, em roxo, arqueada sobre o ramo) e uma
capacitância shunt para o núcleo aterrado \(\Cg\) (em amarelo). A fonte de surto \(\Vsrc\) entra à
esquerda; o neutro (terminal final) é \emph{solidamente aterrado}. É o mesmo circuito da apresentação
animada e da validação cruzada em ATP.}
\label{fig:circuito-pi}
\end{figure}


Do ponto de vista físico, o ramo \(R-L\) armazena energia magnética e dissipa parte dela em resistência, enquanto o capacitor para terra armazena energia elétrica. A interação entre esses armazenamentos faz com que a perturbação de tensão não apareça instantaneamente em toda a bobina: ela se propaga como onda. A impedância de surto aproximada do conjunto é:

\begin{equation}
\Zsurge = \sqrt{\frac{L_{\mathrm{total}}}{C_{\mathrm{total}}}}
       = \sqrt{\frac{0.01}{10^{-9}}}
       \approx \SI{3162}{\ohm}.
\end{equation}

Uma estimativa de tempo de propagação associada ao par \(L_{\mathrm{total}}, C_{\mathrm{total}}\) é:

\begin{equation}
\Ttravel = \sqrt{L_{\mathrm{total}}C_{\mathrm{total}}}
         = \sqrt{0.01\cdot 10^{-9}}
         \approx \SI{3.16}{\micro\second}.
\end{equation}

A saída está em circuito aberto. Para uma onda incidente, essa condição implica corrente terminal nula e tende a produzir reflexão de tensão com mesmo sinal. A duplicação aproximada do pico na extremidade aberta é coerente com os resultados numéricos: a tensão de entrada atinge \SI{1000}{V}, enquanto a saída chega a cerca de \SI{2039}{V}. Esse valor não é uma simples regra ideal, porque a rede discreta, a resistência e a forma temporal do impulso modificam a resposta.

O modelo T, também implementado no projeto, desloca a capacitância de cada seção para um ponto intermediário. Em termos de engenharia, Pi e T são aproximações discretas diferentes de uma mesma estrutura distribuída. A comparação entre eles ajuda a avaliar se a interpretação física depende fortemente da topologia numérica escolhida.
